\documentclass[11pt]{article}
\usepackage[utf8]{inputenc}
\usepackage{amsmath,amssymb,amsthm}
\usepackage{geometry}
\usepackage{hyperref}
\usepackage{mathtools}
\usepackage{xcolor}
\usepackage{listings}
\usepackage{courier}
\geometry{margin=1in}

\lstset{
  basicstyle=\footnotesize\ttfamily,
  breaklines=true,
  columns=fullflexible,
  frame=single,
  numbers=left,
  numberstyle=\tiny,
  captionpos=b,
  xleftmargin=2pt,
  xrightmargin=2pt,
  escapeinside={(*@}{@*)}
}

\begin{document}

\begin{center}
    {\LARGE \textbf{Nondimensionalization of CAR--T model}}\\[6pt]
    \small Copy-paste-ready \LaTeX\ file containing the derivation and Python code.
\end{center}

\bigskip

\section*{Dimensional model (clean / explicit)}
The dimensional model (as in the notebook) is
\begin{align}
\frac{dT}{dt} &= r\,T\!\left(1-\frac{T}{K}\right) - \kappa\frac{C}{C+C_{50}}\,T,
\label{eq:T_dim}\\[6pt]
\frac{dC}{dt} &= p\,C\cdot\frac{T^{h_1}}{T^{h_1}+T_{50,1}^{h_1}}
\;-\; \ell\,C \;-\; \delta_C C \;-\; q\,C\cdot\frac{T/C}{T/C + s},
\label{eq:C_dim}\\[6pt]
\frac{dM}{dt} &= \ell\,C - b\,M\frac{T^{h_2}}{T^{h_2}+T_{50,2}^{h_2}} - \delta_M M.
\label{eq:M_dim}
\end{align}

\medskip

\section*{Choice of characteristic scales}
We choose the characteristic scales as discussed in the meeting and notes:
\[
\text{Tumor scale: } K,\qquad
\text{Effector scale: } C_{50},\qquad
\text{Memory scale: } C_{50},\qquad
\text{Time scale: } \tau = \frac{1}{r}.
\]
Define the dimensionless time and variables
\[
s = r t,\qquad
x(s)=\frac{T(t)}{K},\qquad
y(s)=\frac{C(t)}{C_{50}},\qquad
z(s)=\frac{M(t)}{C_{50}}.
\]
Therefore \(\dfrac{d}{dt} = r\,\dfrac{d}{ds}\).

\bigskip

\section*{Tumor equation (nondimensionalization)}
Starting from \eqref{eq:T_dim}
\[
\frac{dT}{dt} = rT\left(1-\frac{T}{K}\right) - \kappa\frac{C}{C + C_{50}}T.
\]
Substitute \(T=Kx\), \(C=C_{50}y\) and divide by \(rK\). Define
\[
\alpha := \frac{\kappa}{r}.
\]
This yields
\[
\boxed{\displaystyle \frac{dx}{ds} = x(1-x) - \alpha\,\frac{y}{1+y}\,x.}
\]

\bigskip

\section*{Effector (C) equation (nondimensionalization)}
Starting from \eqref{eq:C_dim}
\[
\frac{dC}{dt} = p\,C\cdot\frac{T^{h_1}}{T^{h_1}+T_{50,1}^{h_1}} - \ell\,C - \delta_C C - q\,C\cdot\frac{T/C}{T/C + s}.
\]
Divide by \(r C_{50}\) and substitute the scaled variables.
Define the dimensionless groups
\[
\beta := \frac{p}{r},\quad
\lambda := \frac{\ell}{r},\quad
\delta_c := \frac{\delta_C}{r},\quad
\gamma := \frac{q}{r}.
\]
Also set
\[
\theta_1 := \frac{T_{50,1}}{K}.
\]
Rewrite the exhaustion-like ratio:
\[
\frac{T/C}{T/C + s}
= \frac{Kx/(C_{50}y)}{Kx/(C_{50}y) + s}
= \frac{x}{x + \mu\,y},\qquad
\mu := \frac{s\,C_{50}}{K}.
\]
Therefore the nondimensional effector equation is
\[
\boxed{\displaystyle
\frac{dy}{ds} = \beta\, y\cdot\frac{x^{h_1}}{x^{h_1}+\theta_1^{h_1}} - (\lambda + \delta_c)\,y - \gamma\, y\cdot\frac{x}{x+\mu y}.}
\]

\bigskip

\section*{Memory (M) equation (nondimensionalization)}
Starting from \eqref{eq:M_dim}
\[
\frac{dM}{dt} = \ell\,C - b\,M\frac{T^{h_2}}{T^{h_2}+T_{50,2}^{h_2}} - \delta_M M.
\]
Divide by \(r C_{50}\) and define
\[
b' := \frac{b}{r},\quad \delta_m := \frac{\delta_M}{r},\quad \theta_2 := \frac{T_{50,2}}{K}.
\]
We obtain
\[
\boxed{\displaystyle
\frac{dz}{ds} = \lambda\, y - b'\, z\cdot\frac{x^{h_2}}{x^{h_2} + \theta_2^{h_2}} - \delta_m\, z.}
\]

\bigskip

\section*{Final nondimensional system}
Collecting all three equations:
\[
\boxed{\begin{aligned}
\frac{dx}{ds} &= x(1-x) - \alpha\,\frac{y}{1+y}\,x,\\[6pt]
\frac{dy}{ds} &= \beta\, y\cdot\frac{x^{h_1}}{x^{h_1}+\theta_1^{h_1}} - (\lambda + \delta_c)\,y - \gamma\, y\cdot\frac{x}{x+\mu y},\\[6pt]
\frac{dz}{ds} &= \lambda\, y - b'\, z\cdot\frac{x^{h_2}}{x^{h_2} + \theta_2^{h_2}} - \delta_m\, z.
\end{aligned}}
\]

\subsection*{Dimensionless parameter definitions}
\[
\begin{aligned}
\alpha &= \dfrac{\kappa}{r}, &\quad
\beta &= \dfrac{p}{r}, &\quad
\lambda &= \dfrac{\ell}{r},\\[4pt]
\delta_c &= \dfrac{\delta_C}{r}, &\quad
\delta_m &= \dfrac{\delta_M}{r}, &\quad
\gamma &= \dfrac{q}{r},\\[4pt]
b' &= \dfrac{b}{r}, &\quad
\theta_i &= \dfrac{T_{50,i}}{K}\ (i=1,2), &\quad
\mu &= \dfrac{s\,C_{50}}{K}.
\end{aligned}
\]

\bigskip

\section*{Assumptions / notes}
\begin{itemize}
  \item Time was nondimensionalized with \(\tau = 1/r\). Choosing another time scale (e.g. \(1/p\)) is straightforward.
  \item Both \(C\) and \(M\) were scaled by \(C_{50}\) for simplicity.
  \item Hill functions are dimensionless because \(T\) was normalized by \(K\).
  \item The exhaustion-like term was rewritten as \(\dfrac{x}{x+\mu y}\). If the original form differs, change the algebra accordingly.
\end{itemize}

\end{document}
